\documentclass[../main.tex]{subfiles}
% Ce marqueur n'existe que lorsque ce fichier est compilé seul.
\newcommand{\SubfileStandaloneFlag}{}

\begin{document}

% En compilation individuelle, la mention légale est affichée.
% Dans le livre complet, elle figure déjà sur la page de garde.
\ifdefined\SubfileStandaloneFlag
  \PageLegaleChapitre
\fi

\chapter{Équations et inéquations du 1er degré}
\label{chap:equations-inequations}

% =====================================================
\section{Les équations du 1er degré à une inconnue}
% =====================================================

\subsection{Définition et vocabulaire}

Une \textbf{équation} est une égalité mathématique qui contient une ou plusieurs lettres appelées \textbf{inconnues} (souvent notées $x$). Résoudre une équation, c'est trouver la ou les valeurs de l'inconnue qui rendent l'égalité vraie. Cette valeur forme l'\textbf{ensemble des solutions} (noté $S$).

\begin{definition}
Une équation du 1er degré à une inconnue peut toujours se ramener à la forme générale :
\[
ax + b = 0 \quad (\text{avec } a \neq 0)
\]
\end{definition}

\subsection{Principes d'équivalence (Règles d'or)}

Pour résoudre une équation, on cherche à isoler l'inconnue $x$ en effectuant des opérations des deux côtés de l'égalité sans en modifier les solutions.

\begin{propriete}[Règles de transformation]
On ne change pas les solutions d'une équation si :
\begin{enumerate}[label=\textbf{(\arabic*)}]
  \item On \textbf{ajoute} ou on \textbf{soustrait} un même nombre (ou une même expression) aux deux membres.
  \item On \textbf{multiplie} ou on \textbf{divise} les deux membres par un même nombre \textbf{non nul}.
\end{enumerate}
\end{propriete}

\subsection{Méthode de résolution pas à pas}

\begin{methode}
Pour résoudre une équation type $3x - 5 = x + 7$ :
\begin{enumerate}
  \item \textbf{Transposer les termes en $x$ d'un côté} (souvent à gauche) et les termes constants de l'autre (en changeant leur signe lors du passage du "pont" $=$).
  \item \textbf{Réduire} chaque membre ($3x - x = 2x$ et $7 + 5 = 12$).
  \item \textbf{Isoler $x$} en divisant par le coefficient qui le multiplie ($2x = 12 \implies x = \frac{12}{2} = 6$).
  \item \textbf{Conclure} avec l'ensemble des solutions : $S = \{6\}$.
\end{enumerate}
\end{methode}

\begin{attention}
Lorsqu'il y a des parenthèses précédées d'un signe « $-$ », n'oublie pas de \textbf{changer tous les signes} à l'intérieur lors de la suppression des parenthèses : $-(2x - 3) = -2x + 3$.
\end{attention}

% =====================================================
\section{Les inéquations du 1er degré}
% =====================================================

\subsection{Définition}

Une \textbf{inéquation} ressemble à une équation, mais au lieu d'un signe « $=$ », elle utilise l'un des quatre symboles d'inégalité : $<$ ( strictement inférieur ), $\leq$ (inférieur ou égal), $>$ (strictement supérieur) ou $\geq$ (supérieur ou égal).

\begin{definition}
Résoudre une inéquation, c'est trouver tous les nombres qui vérifient l'inégalité. La solution n'est généralement plus un nombre isolé, mais un \textbf{intervalle de nombres} (ou une demi-droite graduée).
\end{definition}

\subsection{La règle d'or des inéquations}

Les principes de résolution sont les mêmes que pour les équations, à une exception fondamentale près :

\begin{attention}
Si l'on multiplie ou si l'on divise les deux membres d'une inéquation par un \textbf{nombre négatif}, il faut \textbf{impérativement inverser le sens du signe d'inégalité} ($<$ devient $>$, $\leq$ devient $\geq$, etc.).
\end{attention}

\textbf{Exemple :} 
$-2x > 6 \quad\Longrightarrow\quad x < \frac{6}{-2} \quad\Longrightarrow\quad x < -3$. 
(Le signe $>$ s'est transformé en $<$ car on a divisé par $-2$).

% =====================================================
\section{Récapitulatif}
% =====================================================

\begin{tcolorbox}[colback=gray!6,colframe=gray!60!black,title={\bfseries À retenir}]
\textbf{Équations} ($=$) : 
Passage d'un terme de l'autre côté $\rightarrow$ on change son signe ($+$ devient $-$, $\times$ devient $\div$).

\textbf{Inéquations} ($<, \leq, >, \geq$) : 
Même méthode, \textbf{mais} si division/multiplication par un nombre négatif $\rightarrow$ \textbf{on inverse le sens de l'inégalité}.

\textbf{Représentation graphique} (pour les inéquations) :
\begin{itemize}
  \item Crochet fermé $[\ ]$ si l'inégalité est large ($\leq$ ou $\geq$).
  \item Crochet ouvert $]\ [$ si l'inégalité est stricte ($<$ ou $>$).
  \item La flèche indique la direction des solutions sur la droite graduée.
\end{itemize}
\end{tcolorbox}

\newpage
% =====================================================
\section{Exercices}
% =====================================================

\subsection{Série A -- Résolution d'équations}

\begin{enumerate}[label=\textbf{\arabic*.}]

\item Résous les équations suivantes dans $\mathbb{R}$ :
\begin{enumerate}[label=\alph*)]
  \item $4x - 5 = 19$ \qquad
  \item $7x + 3 = 2x - 7$ \qquad
  \item $5(x - 2) = 3x + 4$ \qquad
  \item $3(2x - 1) - 2(x + 4) = 5$
\end{enumerate}

\item Résous les équations avec des fractions :
\begin{enumerate}[label=\alph*)]
  \item $\dfrac{x}{3} + \dfrac{x}{2} = 5$ \qquad
  \item $\dfrac{2x - 1}{4} = \dfrac{x + 3}{2}$
\end{enumerate}

\end{enumerate}

\subsection{Série B -- Résolution d'inéquations}

\begin{enumerate}[label=\textbf{\arabic*.},start=3]

\item Résous les inéquations suivantes, donne l'ensemble des solutions $S$ et représente-les sur une droite graduée :
\begin{enumerate}[label=\alph*)]
  \item $3x + 7 \le 22$ \qquad
  \item $5 - 2x > 11$ \qquad
  \item $4x - 3 \ge 7x + 9$ \qquad
  \item $2(x - 3) < 5(x + 1) - 4$
\end{enumerate}

\end{enumerate}

\subsection{Série C -- Problèmes de mise en équation}

\begin{enumerate}[label=\textbf{\arabic*.},start=4]

\item \textbf{L'âge de la troupe.} Un chef scout a trois fois l'âge de son économe. Dans 10 ans, il n'aura plus que deux fois son âge. Quel est l'âge de chacun aujourd'hui ?

\item \textbf{Le parcours de rando.} Pour organiser un grand jeu en forêt, la patrouille dispose d'un budget de $50\ \text{\euro}$. La location du matériel de base coûte $20\ \text{\euro}$, et chaque kit de balisage supplémentaire coûte $4\ \text{\euro}$. 
\begin{enumerate}[label=\alph*)]
  \item Écris une inéquation traduisant le fait que le coût total ne doit pas dépasser le budget.
  \item Résous cette inéquation pour trouver le nombre maximum de kits de balisage que la patrouille peut s'offrir.
\end{enumerate}

\end{enumerate}

\subsection{Défi bonus}

Résous l'équation paramétrique suivante d'inconnue $x$ (où $m$ est un nombre réel non nul) :
\[
m(x - 1) = 2x + 3
\]
Exprime $x$ en fonction de $m$ et discute selon les valeurs de $m$.

\newpage
% =====================================================
\section{Corrigés}
% =====================================================

\subsection*{Série A}

\textbf{1.} 
\begin{enumerate}[label=\alph*)]
  \item $4x = 24 \implies x = 6 \quad \Rightarrow S = \{6\}$
  \item $5x = -10 \implies x = -2 \quad \Rightarrow S = \{-2\}$
  \item $5x - 10 = 3x + 4 \implies 2x = 14 \implies x = 7 \quad \Rightarrow S = \{7\}$
  \item $6x - 3 - 2x - 8 = 5 \implies 4x - 11 = 5 \implies 4x = 16 \implies x = 4 \quad \Rightarrow S = \{4\}$
\end{enumerate}

\textbf{2.}
\begin{enumerate}[label=\alph*)]
  \item On réduit au même dénominateur ($6$) : $\dfrac{2x}{6} + \dfrac{3x}{6} = \dfrac{30}{6} \implies 5x = 30 \implies x = 6 \quad \Rightarrow S = \{6\}$
  \item On multiplie par $4$ : $2x - 1 = 2(x + 3) \implies 2x - 1 = 2x + 6 \implies -1 = 6$ (impossible) $\quad \Rightarrow S = \emptyset$
\end{enumerate}

\subsection*{Série B}

\textbf{3.}
\begin{enumerate}[label=\alph*)]
  \item $3x \le 15 \implies x \le 5 \quad \Rightarrow S = ]-\infty, 5]$
  \item $-2x > 6 \implies x < -3$ (on a inversé le sens car division par $-2$) $\quad \Rightarrow S = ]-\infty, -3[$
  \item $-3x \ge 12 \implies x \le -4$ (on a inversé le sens car division par $-3$) $\quad \Rightarrow S = ]-\infty, -4]$
  \item $2x - 6 < 5x + 5 - 4 \implies -3x < 7 \implies x > -\dfrac{7}{3}$ (on a inversé le sens) $\quad \Rightarrow S = ]-\dfrac{7}{3}, +\infty[$
\end{enumerate}

\subsection*{Série C}

\textbf{4.} Soit $x$ l'âge de l'économe aujourd'hui. Le chef a $3x$. 
Dans 10 ans : l'économe aura $x + 10$ et le chef aura $3x + 10$.
Équation : $3x + 10 = 2(x + 10) \implies 3x + 10 = 2x + 20 \implies x = 10$.
\textit{Conclusion :} L'économe a $10$ ans et le chef a $30$ ans.

\textbf{5.} Soit $x$ le nombre de kits de balisage.
a) $20 + 4x \le 50$ \quad 
b) $4x \le 30 \implies x \le 7\text{,}5$. 
Comme on ne peut acheter qu'un nombre entier de kits, la patrouille peut s'offrir au maximum \textbf{7 kits}.

\subsection*{Défi bonus}

$mx - m = 2x + 3 \implies mx - 2x = m + 3 \implies x(m - 2) = m + 3$.
Si $m \neq 2$, $x = \dfrac{m + 3}{m - 2}$. (Si $m = 2$, l'équation devient $0x = 5$, ce qui est impossible).

\end{document}
